\documentclass[12pt, letterpaper]{article}

\usepackage[margin=1in]{geometry}
\usepackage{amsmath, amssymb}
\usepackage{booktabs}
\usepackage{array}
\usepackage{hyperref}
\usepackage{enumitem}
\usepackage{parskip}
\usepackage{graphicx}
\usepackage{caption}
\usepackage{algorithm}
\usepackage{algpseudocode}
\usepackage{multirow}
\usepackage{setspace}

\hypersetup{colorlinks=false, pdfborder={0 0 0}}

\newcolumntype{L}[1]{>{\raggedright\arraybackslash}p{#1}}
\newcolumntype{C}[1]{>{\centering\arraybackslash}p{#1}}

\pagestyle{plain}

\begin{document}

% ── Cover Page ───────────────────────────────────────────────────────────────
\begin{titlepage}
    \centering
    \vspace*{2cm}
    {\large SYSC 5108 W}\\[0.3cm]
    {\large Deep Learning}\\[1.5cm]
    {\LARGE \textbf{Land Use and Land Cover Classification\\[0.3cm]
    from Satellite Imagery Using\\[0.3cm]
    Convolutional and Transformer-Based\\[0.3cm]
    Deep Learning Models}}\\[2cm]
    {\large Tan Tran}\\[0.2cm]
    {\normalsize Student ID: 101058868}\\[0.3cm]
    {\normalsize Carleton University}\\[2cm]
    {\normalsize April 2026}
\end{titlepage}

% ── Abstract ─────────────────────────────────────────────────────────────────
\newpage
\section*{Abstract}
\addcontentsline{toc}{section}{Abstract}

This report presents the development and evaluation of three deep learning models for
land use and land cover (LULC) classification from satellite imagery using the EuroSAT
benchmark dataset. The work is motivated by satellite-based ghost gear detection efforts
at Fisheries and Oceans Canada (DFO), where proprietary operational data cannot be used
in an academic setting. EuroSAT, derived from Sentinel-2 imagery provided by the
European Space Agency, serves as the public proxy dataset for developing and validating
the proposed methods.

Three models are trained and evaluated on a stratified 70/15/15 train/validation/test
split of the 27,000-image, 10-class dataset: a Baseline CNN trained from scratch
achieving 93.26\% test accuracy, a ResNet50 with two-phase ImageNet transfer learning
achieving 97.51\%, and a Vision Transformer ViT-B/16 with ImageNet-21k transfer learning
achieving 98.77\%. The ViT-B/16 outperforms both convolutional baselines and exceeds the
original EuroSAT benchmark of 98.57\% reported by Helber et al.\ using ResNet50.

Confusion matrix analysis reveals that vegetation-type classes --- particularly
Permanent Crop, Herbaceous Vegetation, and Pasture --- present the greatest
classification difficulty across all models, accounting for the majority of
misclassifications. River confusion with Highway was also identified as a persistent
and unexpected error pattern. Results confirm that transfer learning provides a
a clear improvement over training from scratch, and that attention-based Transformer
architectures outperform convolutional models for satellite scene classification when
sufficiently pre-trained.

\newpage

% ── Table of Acronyms ─────────────────────────────────────────────────────────
\section*{Table of Acronyms}
\addcontentsline{toc}{section}{Table of Acronyms}

\begin{tabular}{@{}L{3cm} L{10cm}@{}}
    \toprule
    \textbf{Acronym} & \textbf{Definition} \\
    \midrule
    BN    & Batch Normalisation \\
    CNN   & Convolutional Neural Network \\
    DFO   & Fisheries and Oceans Canada \\
    ESA   & European Space Agency \\
    FC    & Fully Connected \\
    F1    & F1-Score (Harmonic Mean of Precision and Recall) \\
    GPU   & Graphics Processing Unit \\
    LULC  & Land Use and Land Cover \\
    MHSA  & Multi-Head Self-Attention \\
    MLP   & Multi-Layer Perceptron \\
    ReLU  & Rectified Linear Unit \\
    RGB   & Red, Green, Blue \\
    SGD   & Stochastic Gradient Descent \\
    ViT   & Vision Transformer \\
    \bottomrule
\end{tabular}

\newpage

% ── Table of Math Notations ───────────────────────────────────────────────────
\section*{Table of Mathematical Notations}
\addcontentsline{toc}{section}{Table of Mathematical Notations}

\begin{tabular}{@{}L{3.5cm} L{9.5cm}@{}}
    \toprule
    \textbf{Symbol} & \textbf{Description} \\
    \midrule
    $x \in \mathbb{R}^{H \times W \times C}$ & Input image of height $H$, width $W$, $C$ channels \\
    $\hat{y}$           & Predicted class label \\
    $y$                 & True class label \\
    $W_l$               & Weight matrix of layer $l$ \\
    $b_l$               & Bias vector of layer $l$ \\
    $h_l$               & Output (activation) of layer $l$ \\
    $g(\cdot)$          & Activation function \\
    $\mathcal{F}(\cdot)$& Residual mapping function \\
    $\mathcal{L}$       & Cross-entropy loss function \\
    $\eta$              & Learning rate \\
    $\lambda$           & Weight decay coefficient \\
    $\nabla \mathcal{L}$& Gradient of the loss function \\
    $\mu$               & Mean (batch normalisation or channel statistics) \\
    $\sigma$            & Standard deviation \\
    $\epsilon$          & Small constant for numerical stability \\
    $\gamma, \beta$     & Learnable scale and shift parameters (batch normalisation) \\
    $N$                 & Number of image patches (ViT) \\
    $P$                 & Patch size in pixels (ViT) \\
    $D$                 & Embedding dimension (ViT) \\
    $Q, K, V$           & Query, Key, Value matrices (self-attention) \\
    $d_k$               & Dimension of key vectors \\
    $\mathbf{E}$        & Patch projection matrix (ViT) \\
    $\mathbf{E}_{\text{pos}}$ & Positional embedding matrix (ViT) \\
    $M$                 & Total number of training examples \\
    $K$                 & Total number of classes \\
    $p$                 & Dropout probability \\
    \bottomrule
\end{tabular}

\newpage

% ── Table of Contents ─────────────────────────────────────────────────────────
\tableofcontents
\newpage

% ════════════════════════════════════════════════════════════════════════════
\section{Introduction}

\subsection{Motivation}

Fisheries and Oceans Canada (DFO) is the federal department responsible for managing and
protecting Canada's marine and freshwater resources. Among its ongoing mandates is the
detection and removal of ghost gear --- abandoned, lost, or discarded fishing equipment
that poses significant ecological risks to marine habitats and species.

\subsubsection*{The Ghost Gear Problem}

Ghost gear is widely recognised as the deadliest form of marine plastic debris.
Between 600,000 and 800,000 metric tonnes of ghost gear enter the world's oceans
annually \cite{canadiangeographic2021}, and ghost gear accounts for approximately
70\% of all macro-plastics in the ocean by weight \cite{ifaw2022}. Once lost, gear
continues to fish indefinitely in a destructive cycle: lost traps and nets retain
bait that attracts marine species, which become entrapped and in turn attract further
species, repeating until the gear degrades or is retrieved. Animals caught in ghost
gear frequently die through suffocation or exhaustion. Ghost gear also damages
critical marine habitats including coral reefs and mangroves, and poses hazards to
navigation.

The economic consequences are equally significant. Approximately 90\% of species
caught in ghost gear are of commercial value \cite{ifaw2022}. A 2021 study found
that lobster traps lost in Southwest Nova Scotia alone cost harvesters \$175,000 in
annual harvest \cite{canadiangeographic2021}. The cost of traditional mapping efforts
to locate ghost gear in three areas of the Canadian Arctic was estimated at \$635,000
for a single operation \cite{dfopsdc2023}. An AI-based approach, by comparison, has
been estimated to cost at most \$500,000 annually for year-round regional-scale
monitoring, while providing more comprehensive and timely coverage.

\subsubsection*{Canada's Policy Context}

Canada has committed to recovering 100\% of all plastics by 2040 and has taken
concrete steps toward this goal through the Oceans Protection Plan, Ocean Plastics
Charter, Ghost Gear Fund, and Blue Economy Strategy \cite{dfopsdc2023}. The Ghost
Gear Fund supports harvesters and communities in retrieving lost gear and returning
it to the circular economy. Extreme weather events such as Hurricane Fiona have
further increased the rate of gear loss, adding urgency to these efforts.

DFO has identified AI-assisted satellite imagery analysis as a key tool for mapping
ghost gear at operational scale. In 2023, DFO submitted a Results Fund Proposal in
partnership with the Ghost Gear Program to pilot AI solutions for detecting marine
litter and predicting gear accumulation densities \cite{dfopsdc2023}. Marina Petrovic,
Assistant Director of Fisheries Resource Management at DFO, noted that Canada is
in a position to build on existing gear-loss datasets by incorporating AI to model
potential areas of gear accumulation and to assist in interpreting sonar mapping
data, which would allow retrieval efforts to be directed toward areas of highest
accumulation or greatest ecological sensitivity \cite{dfopsdc2023}.

\subsubsection*{Scope of This Project}

DFO currently holds proprietary satellite datasets relevant to this monitoring task;
however, this data is subject to confidentiality restrictions that prevent its use in
an academic project. To develop and validate deep learning classification methods
applicable to satellite imagery, this project uses EuroSAT \cite{helber2019eurosat},
a publicly available benchmark derived from Sentinel-2 imagery provided by the European
Space Agency (ESA). EuroSAT provides a well-established testbed of 27,000 labelled
satellite image patches across 10 land use and land cover classes. The underlying
task --- classifying satellite image patches into distinct categories using deep
learning --- applies to the scene-level analysis required for ghost gear detection.
Methods validated here can serve as a basis for future application to DFO's
restricted operational data.

\subsection{Objectives}

This project trains and evaluates three deep learning models for LULC classification
on EuroSAT:

\begin{enumerate}
    \item \textbf{Baseline CNN} --- a compact convolutional network trained from
          scratch, establishing a lower-bound reference.
    \item \textbf{ResNet50 with Transfer Learning} --- a 50-layer residual network
          pre-trained on ImageNet \cite{he2016resnet}, fine-tuned using a two-phase
          strategy.
    \item \textbf{Vision Transformer ViT-B/16 with Transfer Learning} --- a
          Transformer-based architecture pre-trained on ImageNet-21k
          \cite{dosovitskiy2021vit}, fine-tuned on EuroSAT.
\end{enumerate}

The comparative analysis addresses three questions: (i) how much does pre-training
improve performance on satellite imagery relative to training from scratch, (ii) do
attention-based Transformer architectures outperform convolutional architectures for
LULC classification, and (iii) which land use categories are most susceptible to
misclassification across all models.

% ════════════════════════════════════════════════════════════════════════════
\section{Background}

\subsection{Related Concepts}

\subsubsection{Deep Neural Networks}

A neural network consists of a series of interconnected layers of neurons. Each neuron
accumulates a weighted sum of its inputs from the previous layer, adds a bias, and
applies a non-linear activation function to produce its output:

\begin{equation}
    h_l = g(W_l h_{l-1} + b_l)
    \label{eq:neuron}
\end{equation}

where $W_l$ are the layer weights, $b_l$ is the bias, and $g(\cdot)$ is the activation
function. The most common modern activation function is the Rectified Linear Unit (ReLU):

\begin{equation}
    g(z) = \max(0, z)
    \label{eq:relu}
\end{equation}

Training adjusts weight values by propagating prediction errors backwards through the
network --- a process known as backpropagation --- to minimise a loss function.

\subsubsection{Convolutional Neural Networks}

A Convolutional Neural Network (CNN) applies learnable filters across the input using
a sliding window, producing feature maps that detect local spatial patterns. Batch
normalisation \cite{ioffe2015batch} is applied after convolution to stabilise training:

\begin{equation}
    \hat{h} = \frac{h - \mu_{\text{batch}}}{\sqrt{\sigma_{\text{batch}}^2 + \epsilon}},
    \qquad
    \hat{h}_{\text{scale}} = \gamma \hat{h} + \beta
    \label{eq:bn}
\end{equation}

A max pooling operation follows to reduce spatial dimensions and provide translation
invariance. Deep CNNs stack many such layers, building increasingly abstract feature
representations from low-level edges to complex class-specific patterns.

\subsubsection{Residual Networks}

Training very deep CNNs is hindered by the vanishing gradient problem. ResNets
\cite{he2016resnet} address this through skip connections:

\begin{equation}
    \mathbf{y} = \mathcal{F}(\mathbf{x},\, \{W_i\}) + \mathbf{x}
    \label{eq:resnet}
\end{equation}

where $\mathcal{F}$ is the residual mapping and $\mathbf{x}$ is the identity shortcut.
This gives a direct gradient path during backpropagation, allowing stable training
of networks with tens to hundreds of layers.

\subsubsection{Vision Transformers}

The Vision Transformer (ViT) \cite{dosovitskiy2021vit} reframes image classification
as a sequence modelling problem. An input image is partitioned into $N = HW/P^2$
non-overlapping patches, each linearly projected into a $D$-dimensional embedding.
A learnable class token is prepended and positional embeddings are added:

\begin{equation}
    \mathbf{z}_0 = \bigl[\mathbf{x}_{\text{class}};\;
    \mathbf{x}_p^1 \mathbf{E};\; \ldots;\; \mathbf{x}_p^N \mathbf{E}\bigr]
    + \mathbf{E}_{\text{pos}}
    \label{eq:vit_input}
\end{equation}

The sequence passes through Transformer encoder layers applying Multi-Head
Self-Attention:

\begin{equation}
    \text{Attention}(Q, K, V) =
    \text{softmax}\!\left(\frac{QK^\top}{\sqrt{d_k}}\right) V
    \label{eq:attention}
\end{equation}

The final class token representation is passed to a linear classification head.
Unlike CNNs, self-attention operates globally from the first layer, allowing the model
to reason about relationships between any two image patches regardless of spatial
distance.

\subsubsection{Transfer Learning}

Transfer learning initialises a model with weights learned on a large source dataset
and fine-tunes on the smaller target dataset. The pre-trained weights encode general
visual features that transfer across domains, reducing the amount of target-domain
data and training time required. Yosinski et al.\ \cite{yosinski2014transferable}
showed that lower network layers learn highly transferable features while
upper layers are more task-specific and benefit most from fine-tuning. This motivates
the standard practice of replacing the final classification layer for the target
number of classes and fine-tuning the network on the new dataset.

\subsubsection{Cross-Entropy Loss}

All models are trained using cross-entropy loss:

\begin{equation}
    \mathcal{L}(W) = -\frac{1}{M} \sum_{i=1}^{M} \sum_{k=1}^{K}
    \mathbf{1}[y^{(i)} = k] \log \hat{p}_k^{(i)}
    \label{eq:ce}
\end{equation}

where $M$ is the number of training examples, $K$ is the number of classes, and
$\hat{p}_k^{(i)}$ is the predicted probability for class $k$.

\subsection{Related Work}

Helber et al.\ \cite{helber2019eurosat} introduced EuroSAT and established a benchmark
of 98.57\% accuracy using ResNet50, showing that ImageNet pre-trained CNNs
transfer well to satellite imagery despite the domain shift from natural images.
Marmanis et al.\ \cite{marmanis2016deep} applied ImageNet-pre-trained CNN features to
aerial scene classification and showed accuracy gains over models trained
from scratch. Cheng et al.\ \cite{cheng2017remote} confirmed this across multiple remote
sensing benchmarks, establishing transfer learning as the standard approach for satellite
image classification with limited labelled data.

More recently, Vision Transformers have been applied to remote sensing tasks.
Dosovitskiy et al.\ \cite{dosovitskiy2021vit} showed that ViT models pre-trained on
large datasets achieve strong performance on image classification benchmarks, with
performance generally matching or exceeding CNNs when sufficient pre-training data is
available.

\subsection{Rationale for Chosen Methods}

\subsubsection{Why Deep Learning?}

Satellite images are high-dimensional data with complex spatial relationships that
traditional machine learning approaches with hand-crafted features cannot capture
reliably. Deep learning models learn hierarchical feature representations directly
from raw pixels, giving better accuracy and generalisation across diverse land
cover types.

\subsubsection{Why These Three Architectures?}

The three models are chosen to form a structured comparison. The Baseline CNN
establishes what is achievable from EuroSAT alone. ResNet50 represents the established
CNN-based transfer learning approach and provides a direct comparison with the Helber
et al.\ benchmark. ViT-B/16 tests whether global attention-based reasoning provides
advantages over local convolutional feature extraction for satellite scene
classification, where class identity often depends on spatial layout across the whole
image rather than any single local region.

% ════════════════════════════════════════════════════════════════════════════
\section{System Description}

\subsection{Dataset Description}

The EuroSAT dataset \cite{helber2019eurosat} is derived from Sentinel-2 satellite
imagery provided by the ESA. It contains 27,000 labelled $64 \times 64$ pixel RGB
image patches across 10 LULC classes, available at
\url{https://zenodo.org/records/7711810}.

\begin{table}[h!]
    \centering
    \caption{EuroSAT Class Distribution}
    \label{tab:classes}
    \begin{tabular}{@{}clcc@{}}
        \toprule
        \textbf{ID} & \textbf{Class} &
        \textbf{Total Samples} & \textbf{Test Samples} \\
        \midrule
        0 & Annual Crop           & 3,000 & 450 \\
        1 & Forest                & 3,000 & 450 \\
        2 & Herbaceous Vegetation & 3,000 & 451 \\
        3 & Highway               & 2,500 & 375 \\
        4 & Industrial            & 2,500 & 375 \\
        5 & Pasture               & 2,000 & 300 \\
        6 & Permanent Crop        & 2,500 & 375 \\
        7 & Residential           & 3,000 & 450 \\
        8 & River                 & 2,500 & 375 \\
        9 & Sea / Lake            & 3,000 & 450 \\
        \midrule
          & \textbf{Total}        & \textbf{27,000} & \textbf{4,051} \\
        \bottomrule
    \end{tabular}
\end{table}

A stratified 70/15/15 split is applied, as summarised in Table~\ref{tab:split}.
Stratification ensures each class is proportionally represented in all three subsets.

\begin{table}[h!]
    \centering
    \caption{Dataset Split}
    \label{tab:split}
    \begin{tabular}{@{}lcc@{}}
        \toprule
        \textbf{Split} & \textbf{Allocation} & \textbf{Samples} \\
        \midrule
        Training   & 70\% & 18,900 \\
        Validation & 15\% & 4,050  \\
        Test       & 15\% & 4,051  \\
        \bottomrule
    \end{tabular}
\end{table}

$k$-fold cross-validation was considered but not used for this project. First, three
architecturally distinct models are trained; with $k = 5$, that requires 15 separate
training runs, which is computationally prohibitive given the cost of fine-tuning
ViT-B/16. Second, EuroSAT contains 27,000 images with 2,000--3,000 samples per class,
a scale at which a single stratified split produces a test set of over 4,000 samples
that is already statistically stable --- the reduction in variance from $k$-fold is
minimal at this dataset size.

Per-channel mean and standard deviation are computed from the training split only and
applied to normalise all three splits. Augmentations applied to the training split only
include random horizontal and vertical flips, random rotation ($\pm 15^\circ$), and
colour jitter (brightness and contrast $\pm 0.2$). These transformations are appropriate
for top-down satellite imagery which has no canonical orientation. For ViT-B/16, images
are resized from $64 \times 64$ to $224 \times 224$ to match the model's required
input resolution. No augmentation is applied to validation or test splits.

\subsection{Neural Network Architectures}

\subsubsection{Baseline CNN}

A four-block CNN trained from scratch. Each block follows:
\[
    \text{Conv2D} \;\rightarrow\; \text{BatchNorm} \;\rightarrow\;
    \text{ReLU} \;\rightarrow\; \text{MaxPool}
\]
with channel widths $\{32, 64, 128, 256\}$ and $3 \times 3$ kernels with padding. An
adaptive average pooling layer reduces spatial dimensions to $4 \times 4$, followed by:
\[
    \text{Flatten} \;\rightarrow\; \text{Linear}(4096 \to 512) \;\rightarrow\;
    \text{ReLU} \;\rightarrow\; \text{Dropout}(p{=}0.4) \;\rightarrow\;
    \text{Linear}(512 \to 10)
\]
The model contains approximately 2.1M parameters and accepts $64 \times 64$ RGB input.

\subsubsection{ResNet50 with Transfer Learning}

ResNet50 \cite{he2016resnet} is a 50-layer residual network pre-trained on ImageNet.
The final fully connected layer --- originally configured for 1,000 ImageNet categories
--- is replaced with a new classification head for 10 EuroSAT classes:
\[
    \text{Linear}(2048 \to 256) \;\rightarrow\; \text{ReLU} \;\rightarrow\;
    \text{Dropout}(p{=}0.3) \;\rightarrow\; \text{Linear}(256 \to 10)
\]
The full model contains approximately 25.6M parameters.

\subsubsection{Vision Transformer ViT-B/16}

ViT-B/16 \cite{dosovitskiy2021vit} divides each $224 \times 224$ input into
$N = (224/16)^2 = 196$ patches of $16 \times 16$ pixels, projected into a
$D = 768$-dimensional embedding. The model uses $L = 12$ Transformer encoder layers
with 12 attention heads. The pre-trained classification head is replaced with a linear
layer mapping the 768-dimensional class token to 10 outputs. The full model contains
approximately 86.6M parameters.

\begin{table}[h!]
    \centering
    \caption{Model Architecture Summary}
    \label{tab:arch}
    \begin{tabular}{@{}L{3.2cm} C{2.5cm} C{3.2cm} C{2.5cm} C{2cm}@{}}
        \toprule
        \textbf{Model} & \textbf{Pre-trained} & \textbf{Architecture} &
        \textbf{Parameters} & \textbf{Input} \\
        \midrule
        Baseline CNN & No  & Convolutional        & $\sim$2.1M  & $64{\times}64$ \\
        ResNet50     & Yes & Conv.\ + Residual    & $\sim$25.6M & $64{\times}64$ \\
        ViT-B/16     & Yes & Transformer          & $\sim$86.6M & $224{\times}224$ \\
        \bottomrule
    \end{tabular}
\end{table}

\subsection{Training Process}

All models are implemented in PyTorch \cite{pytorch} using the \texttt{timm} library
\cite{timm} for pre-trained checkpoints and trained on a GPU-accelerated environment.

\subsubsection{Optimiser and Loss}

All models use AdamW \cite{loshchilov2019decoupled} with cross-entropy loss. AdamW
decouples weight decay from the adaptive learning rate, which gives more effective L2
regularisation than standard Adam. Training parameters are listed in
Table~\ref{tab:hyperparams}.

\begin{table}[h!]
    \centering
    \caption{Training Hyperparameters}
    \label{tab:hyperparams}
    \begin{tabular}{@{}L{5.5cm} L{6cm}@{}}
        \toprule
        \textbf{Hyperparameter} & \textbf{Value} \\
        \midrule
        Loss Function    & Cross-Entropy \\
        Optimiser        & AdamW \\
        Learning Rate    & $1 \times 10^{-4}$ \\
        Weight Decay     & $1 \times 10^{-4}$ \\
        Batch Size       & 32 \\
        Epochs           & 20 \\
        LR Scheduler     & ReduceLROnPlateau (factor = 0.5, patience = 3) \\
        Gradient Clipping & max norm = 1.0 \\
        Random Seed      & 42 \\
        \bottomrule
    \end{tabular}
\end{table}

\subsubsection{Two-Phase Fine-Tuning for ResNet50}

Phase 1 (5 epochs): the backbone is frozen and only the new classification head is
trained. This prevents large gradient signals from the randomly initialised head from
corrupting pre-trained backbone weights before the head has stabilised. Phase 2
(15 epochs): all layers are unfrozen and fine-tuned end-to-end at $1 \times 10^{-5}$
(one-tenth the Phase 1 rate), making careful adjustments to backbone features rather
than overwriting them.

\subsubsection{Checkpointing}

At each epoch, if validation accuracy improves, model weights are saved to disk.
After training, the best-checkpoint weights are restored for final test evaluation,
so the evaluated model is the one that generalised best to the validation set
rather than the potentially overfitted final-epoch model.

The training and testing pseudocode are given in Algorithms~\ref{alg:train}
and~\ref{alg:test}.

\begin{algorithm}[h!]
\caption{General Training Procedure}
\label{alg:train}
\begin{algorithmic}[1]
\Require Training data $(\mathcal{X}, \mathcal{Y})$, learning rate $\eta$,
         epochs $N$, weight decay $\lambda$
\Ensure Trained model parameters $\theta^*$
\State Initialise model parameters $\theta$; best\_val\_acc $\leftarrow 0$
\For{epoch $= 1$ to $N$}
    \For{each mini-batch $(x, y)$ in training data}
        \State $\hat{y} \leftarrow \text{model}(x;\, \theta)$
        \State $\mathcal{L} \leftarrow \text{CrossEntropy}(\hat{y}, y)$
        \State $\nabla\mathcal{L} \leftarrow \text{backprop}(\mathcal{L})$;
               clip $\|\nabla\mathcal{L}\| \leq 1.0$
        \State $\theta \leftarrow \text{AdamW}(\theta,\, \nabla\mathcal{L},\,
               \eta,\, \lambda)$
    \EndFor
    \For{each mini-batch $(x, y)$ in validation data}
        \State Accumulate validation loss and accuracy (no gradient)
    \EndFor
    \State Update LR scheduler with validation loss
    \If{val\_acc $>$ best\_val\_acc}
        \State best\_val\_acc $\leftarrow$ val\_acc; save $\theta$ to checkpoint
    \EndIf
\EndFor
\State Restore $\theta^*$ from best checkpoint
\end{algorithmic}
\end{algorithm}

\subsection{Testing Process}

After training, the best checkpoint is evaluated once on the held-out test set.
No augmentation is applied during testing.

\begin{algorithm}[h!]
\caption{Testing Procedure}
\label{alg:test}
\begin{algorithmic}[1]
\Require Test data $(\mathcal{X}_{\text{test}}, \mathcal{Y}_{\text{test}})$,
         trained parameters $\theta^*$
\Ensure Test accuracy, macro F1, confusion matrix
\State $\hat{Y} \leftarrow [\,]$, $Y \leftarrow [\,]$
\For{each mini-batch $(x, y)$ in test data}
    \State $\hat{y} \leftarrow \arg\max \text{model}(x;\, \theta^*)$
    \State Append $\hat{y}$ to $\hat{Y}$; append $y$ to $Y$
\EndFor
\State Compute accuracy, per-class precision/recall/F1, macro F1,
       normalised confusion matrix
\end{algorithmic}
\end{algorithm}

% ════════════════════════════════════════════════════════════════════════════
\section{Numerical Results}

\subsection{Performance Metrics}

Table~\ref{tab:results} summarises test set performance across all three models.
ViT-B/16 achieves the highest overall accuracy and macro F1, outperforming both
convolutional baselines. Transfer learning produces a clear improvement over
training from scratch: ResNet50 improves on the Baseline CNN by 4.25 percentage points
and ViT-B/16 improves further by an additional 1.26 points.

\begin{table}[h!]
    \centering
    \caption{Test Set Performance Summary}
    \label{tab:results}
    \begin{tabular}{@{}L{3.5cm} C{3cm} C{3cm} C{3cm}@{}}
        \toprule
        \textbf{Model} & \textbf{Test Accuracy} &
        \textbf{Macro F1} & \textbf{Parameters} \\
        \midrule
        Baseline CNN & 93.26\% & 0.9302 & $\sim$2.1M  \\
        ResNet50     & 97.51\% & 0.9740 & $\sim$25.6M \\
        ViT-B/16     & 98.77\% & 0.9872 & $\sim$86.6M \\
        \bottomrule
    \end{tabular}
\end{table}

Table~\ref{tab:prf} presents macro-averaged precision, recall, and F1 for each model.
All three metrics follow the same ranking: ViT-B/16 $>$ ResNet50 $>$ Baseline CNN,
and precision and recall are closely matched within each model, indicating no systematic
bias toward false positives or false negatives at the macro level.

\begin{table}[h!]
    \centering
    \caption{Macro-Averaged Precision, Recall, and F1-Score}
    \label{tab:prf}
    \begin{tabular}{@{}L{3.5cm} C{3cm} C{3cm} C{3cm}@{}}
        \toprule
        \textbf{Model} & \textbf{Precision} & \textbf{Recall} & \textbf{F1} \\
        \midrule
        Baseline CNN & 0.9308 & 0.9311 & 0.9302 \\
        ResNet50     & 0.9756 & 0.9728 & 0.9740 \\
        ViT-B/16     & 0.9877 & 0.9868 & 0.9872 \\
        \bottomrule
    \end{tabular}
\end{table}

Table~\ref{tab:perclass} shows per-class F1 scores for all three models. The weakest
classes are consistent across all models: Permanent Crop and River for the Baseline CNN,
and Permanent Crop and Pasture for the two transfer learning models. Forest, Industrial,
and Residential are the strongest classes across all models, reflecting their visually
distinctive signatures at the $64 \times 64$ resolution.

\begin{table}[h!]
    \centering
    \caption{Per-Class F1-Score Across All Models}
    \label{tab:perclass}
    \begin{tabular}{@{}L{4cm} C{3cm} C{3cm} C{3cm}@{}}
        \toprule
        \textbf{Class} & \textbf{Baseline CNN} &
        \textbf{ResNet50} & \textbf{ViT-B/16} \\
        \midrule
        Annual Crop           & 0.9211 & 0.9629 & 0.9833 \\
        Forest                & 0.9802 & 0.9922 & 0.9978 \\
        Herbaceous Vegetation & 0.9001 & 0.9577 & 0.9791 \\
        Highway               & 0.9052 & 0.9813 & 0.9920 \\
        Industrial            & 0.9685 & 0.9880 & 0.9881 \\
        Pasture               & 0.9218 & 0.9468 & 0.9780 \\
        Permanent Crop        & 0.8736 & 0.9441 & 0.9665 \\
        Residential           & 0.9801 & 0.9922 & 0.9955 \\
        River                 & 0.8776 & 0.9800 & 0.9947 \\
        Sea / Lake            & 0.9738 & 0.9945 & 0.9967 \\
        \midrule
        \textbf{Macro F1}     & \textbf{0.9302} &
        \textbf{0.9740} & \textbf{0.9872} \\
        \bottomrule
    \end{tabular}
\end{table}

\subsection{Training and Validation Curves}

Figures~\ref{fig:curves_cnn}, \ref{fig:curves_resnet}, and \ref{fig:curves_vit}
show the training and validation loss and accuracy curves for each model.

\begin{figure}[h!]
    \centering
    \includegraphics[width=\textwidth]{BaselineCNN_curves.png}
    \caption{Baseline CNN training and validation loss and accuracy curves.}
    \label{fig:curves_cnn}
\end{figure}

\begin{figure}[h!]
    \centering
    \includegraphics[width=\textwidth]{ResNet50_curves.png}
    \caption{ResNet50 training and validation loss and accuracy curves.
             The vertical dashed line at epoch 5 marks the Phase 1 / Phase 2
             boundary where the backbone is unfrozen for end-to-end fine-tuning.}
    \label{fig:curves_resnet}
\end{figure}

\begin{figure}[h!]
    \centering
    \includegraphics[width=\textwidth]{ViT-B16_curves.png}
    \caption{ViT-B/16 training and validation loss and accuracy curves.}
    \label{fig:curves_vit}
\end{figure}

\subsection{Confusion Matrices}

Figures~\ref{fig:cm_cnn}, \ref{fig:cm_resnet}, and \ref{fig:cm_vit} show the
row-normalised confusion matrices for each model on the test set. Row normalisation
converts counts to per-class recall, making the matrix comparable across classes
with different sample sizes. Values on the diagonal represent the proportion of
correctly classified samples per class; off-diagonal values indicate confusion
between specific class pairs.

\begin{figure}[h!]
    \centering
    \includegraphics[width=0.85\textwidth]{BaselineCNN_confusion_matrix.png}
    \caption{Normalised confusion matrix for the Baseline CNN (93.26\%).}
    \label{fig:cm_cnn}
\end{figure}

\begin{figure}[h!]
    \centering
    \includegraphics[width=0.85\textwidth]{ResNet50_confusion_matrix.png}
    \caption{Normalised confusion matrix for ResNet50 (97.51\%).}
    \label{fig:cm_resnet}
\end{figure}

\begin{figure}[h!]
    \centering
    \includegraphics[width=0.85\textwidth]{ViT-B16_confusion_matrix.png}
    \caption{Normalised confusion matrix for ViT-B/16 (98.77\%).}
    \label{fig:cm_vit}
\end{figure}

Figure~\ref{fig:cm_comparison} presents all three confusion matrices side by side.

\begin{figure}[h!]
    \centering
    \includegraphics[width=\textwidth]{confusion_matrix_comparison.png}
    \caption{Side-by-side normalised confusion matrix comparison across all three
             models. The diagonal darkens and off-diagonal values decrease
             progressively from Baseline CNN to ResNet50 to ViT-B/16, reflecting
             the accuracy improvements shown in Table~\ref{tab:results}.}
    \label{fig:cm_comparison}
\end{figure}

\subsection{Top Misclassification Pairs}

Table~\ref{tab:misclass} lists the ten most common class-pair confusion errors
aggregated across all three models. These represent the pairs that were most
consistently difficult across all architectures, independent of model-specific
behaviour.

\begin{table}[h!]
    \centering
    \caption{Most Common Misclassifications (Aggregated Across All Models)}
    \label{tab:misclass}
    \begin{tabular}{@{}L{5cm} L{5cm} C{3cm}@{}}
        \toprule
        \textbf{True Class} & \textbf{Predicted As} & \textbf{Count} \\
        \midrule
        Permanent Crop        & Herbaceous Vegetation & 43 \\
        River                 & Highway               & 39 \\
        Permanent Crop        & Annual Crop           & 32 \\
        Herbaceous Vegetation & Permanent Crop        & 29 \\
        Pasture               & Herbaceous Vegetation & 23 \\
        Annual Crop           & Permanent Crop        & 21 \\
        Pasture               & Annual Crop           & 16 \\
        Permanent Crop        & Highway               & 12 \\
        Residential           & Industrial            & 12 \\
        River                 & Annual Crop           & 12 \\
        \bottomrule
    \end{tabular}
\end{table}

\subsection{Model Comparison}

Figure~\ref{fig:comparison} presents overall test accuracy and macro F1 for all three
models in a single bar chart for direct visual comparison.

\begin{figure}[h!]
    \centering
    \includegraphics[width=0.85\textwidth]{model_comparison.png}
    \caption{Model comparison --- test accuracy and macro F1 on the EuroSAT test set.}
    \label{fig:comparison}
\end{figure}

\subsection{Results Discussion}

\subsubsection{Training Curve Analysis}

The Baseline CNN curves (Figure~\ref{fig:curves_cnn}) show training and validation
accuracy tracking closely throughout all 20 epochs, reaching a final validation
accuracy of approximately 93\%. The small gap between training and validation curves
indicates that the model generalises well given its limited capacity, but also that
it has reached its performance ceiling --- additional epochs would yield diminishing
returns without architectural changes or more data.

The ResNet50 curves (Figure~\ref{fig:curves_resnet}) display the characteristic
two-phase behaviour of the fine-tuning strategy. During Phase 1 (epochs 1--5, left
of the dashed boundary line), validation accuracy rises from approximately 75\% to
83\% as the classification head is trained against a frozen backbone that immediately
provides rich pre-trained feature representations. At epoch 5, when the backbone is
unfrozen and the lower learning rate is applied for Phase 2, a sharp improvement
occurs and validation accuracy climbs steadily to approximately 97\% by epoch 20.
Training and validation curves remain closely aligned throughout, indicating effective
generalisation.

The ViT-B/16 curves (Figure~\ref{fig:curves_vit}) show a different pattern. Training
loss descends steeply and approaches near-zero by the final epochs, while validation
loss decreases more slowly and with more oscillation. This growing gap between training
and validation loss in the later epochs indicates mild overfitting. Nevertheless, the
validation accuracy stabilises around 98--99\%, showing strong generalisation to
unseen data despite the training/validation loss gap.

\subsubsection{Confusion Matrix Analysis}

The confusion matrices reveal several consistent patterns across all models. Forest,
Industrial, Sea/Lake, and Residential are the easiest classes to classify, with
recall values at or above 0.98 across all three models. These classes have visually
distinctive signatures at $64 \times 64$ resolution --- Forest has dense uniform
green canopy, Sea/Lake has open water with characteristic spectral properties,
Industrial and Residential have distinct structural textures --- that even the
Baseline CNN can identify reliably.

The vegetation-type classes present the greatest difficulty across all models.
Permanent Crop achieves the lowest F1 score in both the Baseline CNN (0.8736) and
ResNet50 (0.9441), and is the second-lowest in ViT-B/16 (0.9665). The confusion
table confirms this, with Permanent Crop being confused with Herbaceous Vegetation
(43 times), Annual Crop (32 times), and vice versa. At the $64 \times 64$ resolution
and in the RGB colour space, these classes share similar ground cover appearance ---
green, textured surfaces without distinctive structural features --- and their
distinguishing characteristics (crop row spacing, canopy structure) are difficult to
resolve at this scale.

The River-to-Highway confusion (39 occurrences) is the most notable unexpected error.
At $64 \times 64$ resolution, narrow rivers and roads can appear as similar linear
features with comparable widths, especially when water is darker or the image is
captured under varying illumination. This confusion is largely resolved by
ResNet50 and ViT-B/16, whose pre-trained features appear to encode contextual cues
(surrounding vegetation patterns, surface texture) that distinguish water bodies
from paved surfaces.

Pasture is consistently confused with Herbaceous Vegetation and Annual Crop, which
is expected as all three classes represent open green areas at different stages of
agricultural use. The distinction between managed pasture and herbaceous vegetation
is subtle even to human observers in small patches.

\subsubsection{Impact of Transfer Learning}

The 4.25 percentage point accuracy improvement from Baseline CNN (93.26\%) to ResNet50
(97.51\%) measures the value of ImageNet pre-training on this domain-shifted
dataset. The improvement is especially pronounced for classes that were weak in the
Baseline CNN --- River improves from F1 0.8776 to 0.9800 and Highway from 0.9052 to
0.9813 --- suggesting that pre-trained features encode texture and structural
information that is difficult to learn from 18,900 satellite patches alone.

\subsubsection{CNN vs.\ Transformer}

ViT-B/16 outperforms ResNet50 by 1.26 percentage points (98.77\% vs.\ 97.51\%),
showing that global attention-based reasoning offers an improvement over local
convolutional feature extraction for satellite scene classification. The improvement
is distributed across nearly all classes, with the most notable gains in Permanent
Crop (+2.24 F1 points) and Pasture (+3.12 F1 points) --- precisely the vegetation
classes where spatial layout context is most useful for disambiguation.

\subsection{Comparison with Related Results}

Table~\ref{tab:related} compares results against published benchmarks.

\begin{table}[h!]
    \centering
    \caption{Comparison with Related Work}
    \label{tab:related}
    \begin{tabular}{@{}L{3cm} C{1.5cm} L{3.5cm} C{2cm} C{2cm}@{}}
        \toprule
        \textbf{Work} & \textbf{Year} & \textbf{Architecture} &
        \textbf{Dataset} & \textbf{Accuracy} \\
        \midrule
        This work & 2026 & Baseline CNN (scratch) & EuroSAT & 93.26\% \\
        This work & 2026 & ResNet50 (transfer)    & EuroSAT & 97.51\% \\
        This work & 2026 & ViT-B/16 (transfer)    & EuroSAT & \textbf{98.77\%} \\
        \cite{helber2019eurosat} & 2019 & ResNet50 & EuroSAT & 98.57\% \\
        \cite{marmanis2016deep}  & 2016 & CNN (transfer) & UC Merced & 94.50\% \\
        \cite{cheng2017remote}   & 2017 & VGG-VD16 (transfer) & AID & 92.07\% \\
        \bottomrule
    \end{tabular}
\end{table}

The ResNet50 result of 97.51\% falls 1.06 points below the 98.57\% benchmark reported
by Helber et al.\ \cite{helber2019eurosat}. This gap is attributable to differences
in training duration, augmentation strategy, and the absence of test-time augmentation
in this implementation. Notably, ViT-B/16 at 98.77\% exceeds the original EuroSAT
ResNet50 benchmark, establishing a new reference result on this dataset using the same
public split methodology.

% ════════════════════════════════════════════════════════════════════════════
\section*{Conclusion}
\addcontentsline{toc}{section}{Conclusion}

This report trained and compared three deep learning models --- a Baseline CNN,
ResNet50 with two-phase transfer learning, and ViT-B/16 with transfer learning ---
for LULC classification on the EuroSAT satellite imagery dataset, motivated by
ghost gear detection work at DFO.

The experiments produced three clear findings. First, transfer learning provides a
substantial benefit on this domain-shifted dataset: ResNet50's 97.51\% accuracy
represents a 4.25 point improvement over the 93.26\% Baseline CNN trained from
scratch on EuroSAT alone, showing that ImageNet pre-trained features generalise
effectively to satellite imagery. Second, the ViT-B/16 Transformer architecture
outperforms the convolutional ResNet50 by 1.26 points, achieving 98.77\% and
exceeding the original EuroSAT benchmark of 98.57\%. This result supports the view
that global attention-based reasoning offers advantages for satellite scene
classification, where class identity often depends on spatial layout rather than
local texture. Third, the most persistent misclassification errors involve
vegetation-type classes --- particularly Permanent Crop, Herbaceous Vegetation, and
Pasture --- whose distinguishing characteristics are difficult to resolve at the
$64 \times 64$ RGB resolution used in this study.

Future work could address these limitations in two directions. Within EuroSAT, the
full 13-band multispectral variant includes near-infrared and shortwave-infrared
bands with information relevant to distinguishing vegetation types
and water bodies, and would likely reduce the vegetation class confusions identified
here. More broadly, the pre-trained models and fine-tuning pipeline developed in
this project apply to DFO's proprietary satellite data for ghost
gear detection, where the trained feature extractors could be further adapted using
a small amount of labelled operational imagery.

% ─────────────────────────────────────────────────────────────────────────────
\newpage
\section*{References}
\addcontentsline{toc}{section}{References}

\begin{thebibliography}{99}

\bibitem{canadiangeographic2021}
Canadian Geographic,
``What Lies Beneath: Ghost Gear in Our Oceans,''
Canadian Geographic Magazine, 2021.
[Online]. Available:
\url{https://canadiangeographic.ca/articles/what-lies-beneath-ghost-gear-in-our-oceans/}

\bibitem{ifaw2022}
Commonwealth Whaling Round Table / International Fund for Animal Welfare (IFAW),
``Entangled: Understanding the Full Scope of Ghost Gear,''
IFAW, 2022.
[Online]. Available:
\url{https://www.ifaw.org/international/projects/entangled-whale-rescue-global}

\bibitem{dfopsdc2023}
DFO PSDC Team,
``Using AI to Detect Ghost Gear,''
Public Service Data Challenge, Fisheries and Oceans Canada, 2023.

\bibitem{helber2019eurosat}
P.\ Helber, B.\ Bischke, A.\ Dengel, and D.\ Borth,
``EuroSAT: A Novel Dataset and Deep Learning Benchmark for Land Use and Land Cover
Classification,''
\textit{IEEE Journal of Selected Topics in Applied Earth Observations and Remote
Sensing}, vol.\ 12, no.\ 7, pp.\ 2217--2226, 2019.

\bibitem{he2016resnet}
K.\ He, X.\ Zhang, S.\ Ren, and J.\ Sun,
``Deep Residual Learning for Image Recognition,''
in \textit{Proceedings of the IEEE Conference on Computer Vision and Pattern
Recognition (CVPR)}, pp.\ 770--778, 2016.

\bibitem{dosovitskiy2021vit}
A.\ Dosovitskiy, L.\ Beyer, A.\ Kolesnikov, D.\ Weissenborn, X.\ Zhai,
T.\ Unterthiner, M.\ Dehghani, M.\ Minderer, G.\ Heigold, S.\ Gelly,
J.\ Uszkoreit, and N.\ Houlsby,
``An Image is Worth 16$\times$16 Words: Transformers for Image Recognition
at Scale,'' in \textit{Proceedings of the International Conference on Learning
Representations (ICLR)}, 2021.

\bibitem{ioffe2015batch}
S.\ Ioffe and C.\ Szegedy,
``Batch Normalization: Accelerating Deep Network Training by Reducing Internal
Covariate Shift,''
in \textit{Proceedings of the 32nd International Conference on Machine Learning
(ICML)}, pp.\ 448--456, 2015.

\bibitem{deng2009imagenet}
J.\ Deng, W.\ Dong, R.\ Socher, L.-J.\ Li, K.\ Li, and L.\ Fei-Fei,
``ImageNet: A Large-Scale Hierarchical Image Database,''
in \textit{Proceedings of the IEEE Conference on Computer Vision and Pattern
Recognition (CVPR)}, pp.\ 248--255, 2009.

\bibitem{loshchilov2019decoupled}
I.\ Loshchilov and F.\ Hutter,
``Decoupled Weight Decay Regularization,''
in \textit{Proceedings of the International Conference on Learning Representations
(ICLR)}, 2019.

\bibitem{yosinski2014transferable}
J.\ Yosinski, J.\ Clune, Y.\ Bengio, and H.\ Lipson,
``How Transferable are Features in Deep Neural Networks?''
in \textit{Advances in Neural Information Processing Systems (NeurIPS)},
vol.\ 27, 2014.

\bibitem{marmanis2016deep}
D.\ Marmanis, M.\ Datcu, T.\ Esch, and U.\ Stilla,
``Deep Learning Earth Observation Classification Using ImageNet Pretrained
Networks,''
\textit{IEEE Geoscience and Remote Sensing Letters},
vol.\ 13, no.\ 1, pp.\ 105--109, 2016.

\bibitem{cheng2017remote}
G.\ Cheng, J.\ Han, and X.\ Lu,
``Remote Sensing Image Scene Classification: Benchmark and State of the Art,''
\textit{Proceedings of the IEEE}, vol.\ 105, no.\ 10, pp.\ 1865--1883, 2017.

\bibitem{pytorch}
A.\ Paszke et al.,
``PyTorch: An Imperative Style, High-Performance Deep Learning Library,''
in \textit{Advances in Neural Information Processing Systems (NeurIPS)},
vol.\ 32, 2019.

\bibitem{timm}
R.\ Wightman,
``PyTorch Image Models,''
GitHub repository, 2019.
[Online]. Available: \url{https://github.com/rwightman/pytorch-image-models}

\end{thebibliography}

\end{document}
